\section{Introduction}

The deployment of autonomous multi-agent systems into consequential environments --- financial clearing, clinical decision support, agricultural resource management, and critical infrastructure orchestration --- has exposed a fundamental structural gap: conventional architectures can describe what agents do, but they cannot structurally guarantee that every agent's authority traces to a human principal. As agent populations grow dynamically, spawn recursively, and make decisions with real-world consequences, the question of who authorized each decision, through what chain of delegation, and with what scope constraints has become the central unsolved problem in agentic infrastructure. We introduce the Recursive Spawning Protocol (RSP): a structural mechanism, built on the \bxthree three-layer architecture, that makes the delegation chain from every active agent to a human principal a runtime artifact --- immutable, verifiable, and architecturally enforced --- rather than a forensic reconstruction attempted after the fact.

% -------------------------------------------------------
\subsection{The Autonomous Drift Problem}
% -------------------------------------------------------

The failure mode RSP targets is \emph{autonomous drift}: the condition in which an agent continues to operate with no traceable authorization path to a human principal. Drift is not a simple system failure --- it is more insidious. The agent continues to function, to spawn children, to modify state, to make decisions, with no external indication that its authority has become disconnected from any accountable principal. The chain of delegation that originally authorized its provisioning has been severed, obscured, or was never structurally established --- yet the agent persists in operational state.

Formally, we define an agent as \emph{drifted} when its delegation chain $C(n) = \langle n_0, n_1, \ldots, n \rangle$ does not contain a human principal at its origin $n_0$. Equivalently: there exists no human anchor $h(n_0)$ such that the chain terminates at $n_0$ with $h(n_0)$ as the accountable principal. An agent may be drifted with or without being orphaned --- it may have a living parent, but that parent's authority does not trace to any human decision. An agent is \emph{orphaned} when its parent node $\alpha(n)$ has either terminated or become unreachable, severing the delegation chain at the parent. An agent enters the \emph{drift triad} when it is simultaneously orphaned, drifted, and operating --- it has lost all traceable connection to human accountability yet continues to execute.

The drift problem emerges from three structural conditions endemic to conventional multi-agent runtimes. \emph{Chain opacity} arises when the parent pointer $\alpha(n)$ is maintained as a mutable reference: when a parent terminates or is reparented, the child's authorization trace is retroactively modified --- the child now points to whatever parent it was redirected to, and the original chain is obscured. \emph{Scope mutation} arises when an agent's operating scope $R(n)$ can be re-evaluated and expanded at runtime by the agent itself or by a supervisory node --- without Purpose Layer authorization, without human principal approval, and without any tamper-evident record of the expansion decision. \emph{Termination ambiguity} arises when agent termination is event-driven but not structurally required to propagate to children: a parent may exit cleanly without issuing a termination event, leaving children unaware that their parent has departed.

These three conditions produce four concrete failure modes. \emph{Silent orphaning}: a parent terminates without notifying its child; the child remains operational, orphaned, and drifted. \emph{Scope creep drift}: a node expands its operating scope without Purpose Layer authorization; successive expansions migrate the chain's effective operating envelope far from its human anchor's intent. \emph{Re-parenting drift}: a child agent's parent pointer is redirected by an administrative operation, retroactively obscuring the original authorization trace. \emph{Ghost authority}: a node spawns a child without any authorized delegation chain; the child is provisioned but its warrant traces to no authorized chain.

The consequences are compounding in recursive multi-agent systems, where nodes spawn children, share state, and delegate objectives across a growing chain. A single drifted agent propagates its unauthorized scope to all descendants. The resulting task tree may superficially resemble a correctly authorized decomposition --- correct topology, correct node count --- but its effective operating envelope is unrecognizable as the product of its human anchor's intent.

% -------------------------------------------------------
\subsection{Why Existing Approaches Are Insufficient}
% -------------------------------------------------------

The conventional architectural response to unbounded agent spawning is the \emph{depth limit}: the system enforces a maximum chain depth $D_{\max}$, rejecting any spawn that would exceed this bound. Depth limits are necessary --- they prevent infinite spawn cascades --- but they are structurally insufficient to prevent autonomous drift.

We prove this formally. Consider a chain $C = \langle n_0, n_1, \ldots, n_k \rangle$ where $k < D_{\max}$. The depth constraint is satisfied by construction. Now apply scope creep drift: each node $n_i$ for $i \geq 1$ successively expands its operating scope $R(n_i)$ beyond $R(n_{i-1})$ without Purpose Layer authorization. After $k$ successive expansions, $R(n_k)$ may be entirely disjoint from $R(n_0)$. The chain depth is $k \leq D_{\max}$. The chain has a human anchor at $n_0$. Yet $n_k$ is drifted by definition: its operating scope is not a descendant refinement of its human anchor's intent. A depth limit constrains the \emph{number} of spawn generations; it does not constrain the \emph{scope} of each generation's authorized actions.

The drift prevention requirement is formally a conjunction of two constraints:

\begin{equation}
\text{Drift Prevention} \iff \bigl(\text{depth}(C) \leq D_{\max}\bigr) \ \land \ \bigl(\forall i:\ R(n_{i+1}) \subseteq R(n_i)\bigr)
\label{eq:intro-drift-prevention}
\end{equation}

Depth limits enforce only the first conjunct. A system that enforces $D_{\max}$ but not the scope refinement constraint is immune to unbounded spawn cascades but remains fully vulnerable to scope creep drift. Time-to-live (TTL) constraints suffer from the same insufficiency: TTL constrains an agent's \emph{lifetime}, not its \emph{scope within that lifetime}. An agent with a 60-second TTL and unrestricted scope expansion authority can produce consequences exceeding its human anchor's intent within seconds of spawning.

Logging and audit infrastructure can partially reconstruct the chain after the fact, but retroactive reconstruction is categorically different from runtime enforcement. A system that discovers, after an incident, that its agent was drifted has not prevented drift --- it has documented it. In consequential environments where unauthorized actions may be irreversible, runtime enforcement is not an optional enhancement; it is a structural requirement. The deeper reason depth limits fail is that they address the wrong dimension of the accountability problem: they constrain how many spawn generations can exist; they say nothing about what each generation is authorized to do.

% -------------------------------------------------------
\subsection{Core Insight: Immutable Parent Pointer Chain Makes Drift Structurally Impossible}
% -------------------------------------------------------

The Recursive Spawning Protocol's core insight is that autonomous drift is structurally impossible when the parent pointer $\alpha(n)$ is immutable --- when it is established once at provisioning and no actor, under any circumstances, can modify it thereafter.

An immutable parent pointer eliminates the conditions that produce drift by construction. Chain opacity cannot arise: the child's parent pointer always traces to the node that authorized its provisioning, regardless of subsequent events in the system. The delegation chain is not a mutable reference that can be redirected --- it is a first-class runtime artifact established by the Fact Layer write protocol and hash-chained into the tamper-evident forensic ledger. Scope mutation cannot produce drift without Purpose Layer authorization: the Fact Layer pre-commit hook rejects any spawn event where the proposed child scope is not a descendant refinement of the parent's scope. The subset relationship $R(n_{i+1}) \subseteq R(n_i)$ is not a policy convention --- it is a structural invariant enforced at every spawn event. Orphaning without drift is structurally harmless: when a parent terminates, its children remain traceable to their human anchor through the immutable chain.

The immutable parent pointer is not a storage policy or an access control configuration. It is a write protocol constraint enforced by the Fact Layer at the protocol level --- before any write is executed, not after it is observed. Any attempt to overwrite $\alpha(n)$ after provisioning is rejected regardless of the requestor's identity, privileges, or justification. There is no administrative override; there is no privileged actor capable of modifying the chain after the fact. The immutability is structural, not contractual.

This structural immutability enables runtime verification: at any point during execution, the system can construct the complete delegation chain for any agent $n$, verify that the chain terminates at a human anchor $h(n_0)$, and confirm that each node's operating scope is a descendant refinement of the human anchor's intent. This verification does not depend on the agent's self-reporting, on the completeness of external logging infrastructure, or on the honesty of any machine actor in the chain.

The Recursive Spawning Protocol is the minimal structural mechanism that achieves this guarantee. It adds three constraints to conventional agentic runtimes: an immutable parent pointer established at provisioning and protected by Fact Layer write protocol; a Purpose Layer authorization check at every spawn event requiring scope refinement justification and spawn necessity declaration; and a Fact Layer pre-commit hook enforcing both the depth bound and the scope subset relationship as structural invariants. Together, these constraints make autonomous drift architecturally impossible --- not statistically unlikely, not conventionally prevented, but structurally impossible as a matter of protocol enforcement.

% -------------------------------------------------------
\subsection{Paper Roadmap}
% -------------------------------------------------------

The remainder of this paper is organized as follows. Section~\ref{sec:rs-background} situates RSP within agent lifecycle management, accountability in distributed systems, and establishes the \bxthree\ Framework as the minimal sufficient substrate. Section~\ref{sec:drift} provides the formal characterization of autonomous drift, defines the drift triad and its constituent failure sub-modes, and proves formally that depth limits alone are insufficient. Section~\ref{sec:rs-protocol} specifies the Recursive Spawning Protocol in full: its Purpose Layer validation gate, its Bounds Engine depth management, its Fact Layer immutable pointer and forensic ledger, and its chain verification criteria. Section~\ref{sec:rs-integration} maps each RSP component to its governing \bxthree\ layer and demonstrates the structural necessity of the three-layer mapping. Section~\ref{sec:rs-correctness} formally states and proves the drift prevention, chain integrity, and termination correctness properties. Section~\ref{sec:rs-limitations} examines the structural costs of RSP's guarantees. Section~\ref{sec:rs-conclusion} concludes with a summary of RSP's contribution to the accountability of multi-agent AI systems.
